\documentclass[11pt]{article}
\usepackage[letterpaper,margin=.88in]{geometry}
\usepackage{amsmath,amssymb,amsthm,hyperref}
\setlength{\emergencystretch}{2em}
\hypersetup{hidelinks,pdftitle={A sign-sensitive count for exceptional primes}}
\newcommand{\sfk}{\operatorname{sf}}
\newcommand{\ep}{\varepsilon}
\newcommand{\e}{\mathrm e}
\newcommand{\leg}[2]{\left(\frac{#1}{#2}\right)}
\newtheorem{proposition}{Proposition}
\title{A sign-sensitive count for exceptional primes}
\author{Research notes}
\date{}
\begin{document}
\maketitle
Let $n_p$ be the least quadratic nonresidue at an odd prime $p$. Fix
$c>0$, let $H=\log P$ and $Y=cH^2$, and put
$E(P,Y)=\#\{P<p\le2P:n_p>Y\}$. Under Dirichlet GRH the bound obtained
here is $E(P,Y)\ll_c\sqrt P$. We also test a quadratic large sieve
estimate on a signed sum of characters. These bounds allow a single
exception in arbitrarily large intervals.

\section{The identity across scales and its sign symmetry}
Use the complete statistic $Q_p(c)$ of the preceding note, with
$L=\log p$ and the same height $U=\sqrt L$ at every scale.
Suppose all prime signs up to $\eta L^2$ are positive, and fix
$0<c\le\eta$. Above $n>\eta L^2$ the normalized prime-power weight
at scale $c$ is $\sqrt{cL^2}\,\Lambda(n)\chi_p(n)n^{-3/2}$ and the
Fourier phase is $(cL^2/n)^{it}$. Thus the whole tail Fourier sum is
exactly $\sqrt{c/\eta}(c/\eta)^{it}$ times the tail at $\eta$.
In the notation of (5) of the preceding note, this gives
$B_{\chi,\eta/c}(c)=(c/\eta)B_{\chi,1}(\eta)$. Substitution yields
\[
 \frac{4/3-Q_p(c)}{c}=\frac{4/3-Q_p(\eta)}{\eta}+o_{c,\eta}(1).
 \tag{1}
\]
If all primes up to $\eta L^2$ instead have sign $-1$, the low-prime
amplitude changes sign, its squared norm has the same limit, and it
is still asymptotically orthogonal to the tail. Proper powers cost
$O(\sqrt{U/(cL^2)})=o_c(1)$ in the normalized band norm by the
prime-power estimate of the preceding note. Hence (1) and
$Q_p(c)\ge4/3-c/\eta-o_{c,\eta}(1)$ hold in that case as well.
An eventual all-prime gap $Q_p(c)\le4/3-\delta(c)+o_c(1)$, with
$\delta(c)/c$ unbounded along dyadic $c\downarrow0$, would therefore
prove the little-$o$ statements for both $n_p$ and the least prime
$\ell$ with $(\ell/p)=1$. The count below singles out positive signs.

\section{The exact character expansion}
For large $P$ we have $2\le Y<P$. The least nonresidue is prime by
multiplicativity. Write $M=\prod_{\ell\le Y}\ell$, $r=\pi(Y)$, and
$\Pi=\pi(2P)-\pi(P)$. All products over $\ell$ in this note are over primes.
Then
\[
 E(P,Y)=\sum_{P<p\le2P}\prod_{\ell\le Y}\frac{1+(\ell/p)}2
 =2^{-r}\left(\Pi+\sum_{\substack{d\mid M\\d>1}}S_d\right),
 \qquad S_d=\sum_{P<p\le2P}\leg d p .
 \tag{2}
\]
For squarefree $d>1$ the fundamental discriminant $D_d$ is $d$ if
$d\equiv1\bmod4$, and $4d$ otherwise. At these odd primes the
symbol in (2) equals the primitive quadratic character of conductor
$q_d=D_d$ evaluated at $p$. Thus every nonconstant term is a
nonprincipal Dirichlet character in the prime being counted.
The principal contribution $\Pi/2^r$ tends to zero faster than any
fixed negative power of $P$, since $r\sim cH^2/(2\log H)$.

In particular $E(P,Y)=0$ holds exactly when
\[
 2^{-r}\sum_{d\mid M,\,d>1} S_d < 1-\Pi/2^r .
 \tag{3}
\]
An all-positive prime adds $1-2^{-r}$ to the normalized nonprincipal
sum. This is a one-sided test at the scale of a single prime.
Eventual vanishing for each $c=2^{-j}$ yields $n_p=o(\log^2 p)$
by taking $P$ to be successive powers of two and then letting $j$
increase. Conversely, little-$o$ implies the vanishing since
$\log(2P)/\log P\to1$.

\paragraph{The two pseudosquare problems.}
The test at $2$ forces $p\equiv1$ or $7\bmod 8$. For every odd
$\ell\le Y$ reciprocity gives
\[
 \leg{\ell}{p}=\begin{cases}
 (p/\ell),&p\equiv1\bmod8,\\
 (-p/\ell),&p\equiv7\bmod8 .
 \end{cases}
\]
The first class consists of positive prime $Y$-pseudosquares. In
the second, $-p\equiv1\bmod8$ is locally square at the odd primes.
For the expansion by residue class put $s_1=1$ and $s_7=-1$.
The contribution of a single class $a=1,7$ is
\[
 E_a=2^{1-r}\sum_{d\mid M/2}(s_a/d)
             \sum_{\substack{P<p\le2P\\p\equiv a\ (8)}}(p/d),
 \tag{3a}
\]
with Jacobi symbols and value one at $d=1$. For example in class 7
pairing $d$ with $3d$ for $3\nmid d$ produces the factor $1-(p/3)$
inside the prime sum. This exactly retains primes $p\equiv2\bmod3$
with weight two. The alternating coefficients encode a local test;
their signs alone give no additional estimate for its survivors.
These two sets are precisely the prime solutions of Problems I and II of
Lehmer, Lehmer and Shanks \cite{lls}. They have separate tables of
least prime solutions. For example at cutoff $11$ the least positive
solution is $2641=19\cdot139$, the least positive prime solution is
$2689$, and the least prime in the negative test has absolute value
$479$ \cite[Tables 6.23 and 6.25]{lukes}.

Equivalently, let $m(Y)$ be the smaller of the two least prime absolute
values. The goal is $\log m(Y)/\sqrt Y\to\infty$. Indeed each candidate
has $p>Y$ and $n_p>Y$, and a counterexample sequence provides candidates
at $Y=\eta\log^2 p$. Each fixed set of local tests admits prime solutions
by the Chinese remainder theorem and Dirichlet's theorem.
For the least positive pseudosquare among all nonsquares,
Pomerance and Shparlinski report Weinberger's Dirichlet-GRH lower bound
$\exp(a\sqrt Y)$ with a fixed $a>0$. Their Theorem 1 on distribution
in intervals of length $N$ assumes $N\ge\exp(3Y/\log\log Y)$.
Here $N=P=\exp(\sqrt{Y/c})$ is eventually smaller, and that theorem
concerns positive pseudosquares including composites \cite{ps}.

\section{What termwise GRH gives}
For a fixed nonnegative smooth weight $w$ supported in $(1/2,3)$,
equal to one on $[1,2]$, set
\[
 S_w(f)=\frac1H\sum_{p\ {\rm prime}}(\log p)w(p/P)(f/p)
\]
for positive squarefree $f$. The explicit formula under GRH gives
\[
 S_w(1)\asymp_w P/H,\qquad |S_w(f)|\ll_w
 \frac{\sqrt P}{H}\log(2q_f)\quad (f>1).
 \tag{4}
\]
For completeness, the Mellin transform of $w$ at $1/2+it$ decreases
as $(1+|t|)^{-2}$ or faster. Each nontrivial zero contributes at most
$O_w(\sqrt P(1+|t|)^{-2})$ to the von Mangoldt sum. The unit-interval
zero count gives $O_w(\sqrt P\log(2q_f))$ after summation. The
gamma and trivial-zero terms obey the same bound, for instance by
the compact prime-side explicit formula in \cite[Section 3]{cmqr}.
Removing powers of exponent at least two costs $O_w(\sqrt P)$ by
Chebyshev's bound; division by $H$ gives (4). The principal estimate
uses the ordinary prime number theorem.

For the sharp unweighted interval a usual version is
$S_f\ll\sqrt P(H+\log(2q_f))$ for $f>1$. One can get it by smoothing
the dyadic indicator on relative width $P^{-1/2}$, including the
factor $H/\log n$ to remove the prime weight. The two edge errors
in the von Mangoldt sum are $O(\sqrt P H)$. On the critical line the
Mellin transform of the smoothed weight is
$O((1+|t|)^{-1}(1+|t|/\sqrt P)^{-2})$ by integration by parts.
The zero sum is then $O(\sqrt P H(\log(2q_f)+H))$ before division by
$H$. Since the mean of $\log d$ over $d\mid M$ is $\tfrac12\log M$,
(2) becomes
\[
 E(P,Y)=\Pi/2^r+O\!\left(\sqrt P(H+Y)\right).
 \tag{5}
\]
For an upper bound only, majorize the count before expanding by the
smooth prime weight $\log(p)w(p/P)/H$. Every prime in its support
exceeds $Y$ for large $P$. Equation (4) gives the sharper termwise
full-product bound
\[
 E(P,Y)\ll_w \frac{P}{2^rH}+\frac{\sqrt P(1+Y)}{H}.
 \tag{6}
\]
The error here has order $\sqrt P\log P$ at fixed $c$.

\section{Keeping the signs together}
Let $X=P^a$ with $a>0$ fixed. Use the $m$ primes in $(Y/2,Y]$;
$m\sim Y/(2\log Y)$. Take $k=\lfloor\log X/\log Y\rfloor-1$ and
let $\mathcal D_j$ be the products of $j$ distinct primes in this block,
$j=k,k+1$. All are odd squarefree and at most $X$. Assign weight
$a_d=1/(2|\mathcal D_j|)$ on each of the two disjoint levels and
zero elsewhere. For $A(p)=\sum_d a_d(d/p)$, precisely the property
needed for a majorant is $A(p)=1$ at every candidate. If the entire
block has negative signs, the adjacent degrees give $A(p)=0$. For both
levels, $\log\binom mj=j\log(m/j)+O(j)=(a/2+o_c(1))H$ at fixed $a$,
so $\sum a_d=1$ and
\[
 D_2:=\sum_d a_d^2=P^{-a/2+o_{a,c}(1)}.
 \tag{7}
\]

First apply Heath-Brown's quadratic large sieve \cite[Theorem 1]{hb}
to the signed inner sum as a whole. With $X=P$ it gives
\[
 E(P,Y)\le\sum_{m'\le2P}^{*}
      \left|\sum_{d\le P}^{*}a_d(d/m')\right|^2
 \ll_\ep P^{1+\ep}D_2 \ll_{\ep,c} P^{1/2+\ep}.
 \tag{8}
\]
The stars restrict both indices to odd squarefree integers; the first
sum includes each prime that needs to be excluded. This is a joint
estimate retaining cancellation between the divisor characters, valid
unconditionally. The loss $P^\ep$ prevents it from improving (6) at
fixed $\ep$. Its square-root exponent is consistent with the exceptional-set
estimates discussed in \cite[Introduction]{lw}.

For the sharper conditional estimate take $X=P^3$. Keep the same signed
amplifier, square it and sum over primes with the majorant $w$ above.
On this support every prime factor of a coefficient index is $<p$.
Writing $\sfk(de)=de/\gcd(d,e)^2$ for squarefree $d,e$, this gives
\[
 E(P,Y)\ \le\ D_2 S_w(1)+
    \underbrace{\sum_{d\ne e}a_d a_e S_w(\sfk(de))}_{\mathcal R}.
 \tag{9}
\]
The diagonal is $o(1)$ by (7). For $d\ne e$ the conductor is at most
$4X^2$. Hence (4) and $\sum a_d=1$ give
\[
 |\mathcal R|\ll_w\sqrt P,\qquad E(P,c\log^2 P)\ll_c\sqrt P .
 \tag{10}
\]
This improves the direct smoothed bound (6) by a logarithm. The gain
comes from lowering the conductors in the majorant; (10) uses no
extra cancellation in the off-diagonal sum. A bound
$\mathcal R\le1-\delta_c$ for some $\delta_c>0$ would suffice, since
the diagonal tends to zero. Equation (8) tests actual joint cancellation
and does not provide such a bound. One can also choose a nonnegative
majorant in (9) with extra local tests inserted, but those tests
must then be accounted for in the conductors.

\paragraph{A limitation of this termwise estimate.}
Even optimizing the nonnegative weights of such a divisor amplifier
cannot make its termwise majorant smaller than order $\sqrt P$.
Indeed, consider any probability weights $a_d$ on squarefree
divisors of $M$ and the numerical expression
\[
 \frac{P}{H}D_2+\frac{\sqrt P}{H}
   \sum_{d\ne e}a_d a_e\log(2q_{\sfk(de)}).
 \tag{11}
\]
If $D_2\ge H/\sqrt P$, the first term is at least $\sqrt P$. Otherwise
choose independently $d,e$ with these weights. For fixed $d$ the map
$e\mapsto\sfk(de)$ is injective, since it toggles membership of each
prime divisor of $d$. Thus
\[
 \Pr(\sfk(de)\le P^{1/4}\mid d)
 \le \sqrt{P^{1/4}D_2}\le P^{-1/8}\sqrt H .
\]
With probability $1-o(1)$ the nonprincipal conductor exceeds $P^{1/4}$,
so the second term of (11) is at least $(1/4-o(1))\sqrt P$.
This is a lower bound on the expression used to majorize the signed
sum, not on the count of primes. The unresolved step is control of
the signed average at unit scale in (3), or a comparably strong
one-sided estimate in (9). No bound for the growth of $n_p$ improves here.

\begin{thebibliography}{9}\small
\bibitem{lls} D.~H.~Lehmer, E.~Lehmer and D.~Shanks,
\emph{Integer sequences having prescribed quadratic character},
\href{https://oeis.org/A002189/a002189.pdf}{Problems I and II}.
\bibitem{lukes} R.~F.~Lukes, \emph{A Very Fast Electronic Number Sieve},
\href{https://mspace.lib.umanitoba.ca/server/api/core/bitstreams/611cd835-b7d9-4b12-bc2a-aaf3f0576c04/content}{Chapter 6}.
\bibitem{ps} C.~Pomerance and I.~Shparlinski,
\emph{On pseudosquares and pseudopowers},
\href{https://math.dartmouth.edu/~carlp/PDF/pseudo3.pdf}{Section 1.1 and Theorem 1}.
\bibitem{cmqr} E.~Carneiro, M.~Milinovich, E.~Quesada-Herrera and A.~Ramos,
\emph{Fourier optimization, the least quadratic non-residue, and the least
prime in an arithmetic progression}, \href{https://arxiv.org/html/2404.08380v2#S3}{Section 3}.
\bibitem{hb} D.~R.~Heath-Brown, \emph{A mean value estimate for real
character sums}, \href{https://doi.org/10.4064/aa-72-3-235-275}{Acta Arith. 72 (1995), 235--275}.
\bibitem{lw} Y.-K.~Lau and J.~Wu, \emph{On the least quadratic non-residue},
\href{https://hkumath.hku.hk/~yklau/p/34.pdf}{Introduction}.
\end{thebibliography}
\end{document}
